\section{Final Project Selection}
Having spent the last few weeks researching various topics as outlined in my mind map. I have now decided on my EPQ project.

Considering the topics I researched and identified in my mind map, I want to integrate Digital workflows, , human learning, web development and tools useful for everyday life. Therefore I have decided to investigate building a software Note storage system designed for students to enable them to quickly search their notes and organise their educational and personal lives. While similar tools do exist (and I have compared in my research section) I haven't found any which fit my (and others) needs and don’t have huge sacrifices that must be made.

While my list of requirements is ever changing (and will be outlined throughout my diary as they progress) these are my base musts for my project to be successful:
\begin{itemize}
	\item Ability to view and edit notes from anywhere, and on any device (laptop, phone, pc and even a smart fridge)
	\item Ability to quickly find any notes anywhere and make logical links between notes for exploration
	\item Robust collaboration tools for working with peers and sharing information
	\item Ability to store text, images and even handwriting within notes
	\item Simple setup and file management  
\end{itemize}
My end goal is to produce software I can use for all my studies and for collating all of my EPQ notes (as OneNote doesn’t quite cover my needs - hence the need for improvement in the space) 

Whilst choosing my project I did have some concerns regarding overlap with the A Level computer science course I am taking. I therefore communicated with both my Computer science teacher and Course mentor to gather their opinions. After an in-person conversation we decided that whilst there are similarities to an extent between a potential computer science NEA (which doesn't start until Y13) and my project,  my project would be distinct enough providing I document all of my learnings to great depth and leaned towards the impact research on human learning has had than  towards the direction of my final software artifact

From past experience of both projects I have worked on and have watched from afar that 'scope creep' can (more often than not) have a huge negative impact on the end produced result. To mitigate this in addition to the core values mentioned above I am going to keep an updating document linked to my timeline which included features I know that can be feasibly added in and some stretch goals that I will add only if time permits. I am hopeful that this management and planning will result in me achieving my goals and possibly even developing my project into a tool I can carry with my to university and beyond, but all in good time.

Below is a table of my current project requirements and how I plan to tackle them. This will be updated as I progress through my project and will be used to help me keep track of my progress and ensure I am on track to meet my goals.


\begin{tabularx}{\textwidth}{
    >{\raggedright\arraybackslash}l 
    >{\raggedright\arraybackslash}X 
    >{\raggedright\arraybackslash}X
}
\toprule
\textbf{Task/Feature} & \textbf{Description} & \textbf{Action Plan} \\
\midrule

\rowcolor{gray!10}
\multicolumn{3}{l}{\textbf{High Priority \& High Impact (Do First)}} \\

Core software development & Building the foundational system for note-taking, organization, and accessibility. & Begin development with a focus on core functionalities. \\

Cross-device accessibility & Ensuring the system works on multiple devices (laptop, phone, smart fridge, etc.). & Research frameworks that support multi-platform use. \\

Fast search \& logical links & Implementing a system to quickly find notes and establish relationships between them. & Prioritize efficient indexing and search algorithms. \\

Human learning research integration & Designing the tool based on effective study methods. & Document findings and integrate best practices. \\

\rowcolor{gray!10}
\multicolumn{3}{l}{\textbf{High Priority \& Low Impact (Plan for Later)}} \\

Collaboration tools & Features for peer sharing and real-time collaboration. & Plan architecture but implement later. \\

Handwriting \& multimedia support & Adding handwritten note storage and image support. & Research existing libraries for ease of integration. \\

File management & Ensuring an intuitive and simple organization system. & Define structure before implementation. \\

\rowcolor{gray!10}
\multicolumn{3}{l}{\textbf{Low Priority \& High Impact (Delegate/Streamline)}} \\

Stretch goals \& enhancements & Additional features that improve usability but aren’t core. & Maintain an evolving feature backlog. \\

Advanced AI search suggestions & Smart recommendations for note retrieval. & Implement if time allows. \\

\rowcolor{gray!10}
\multicolumn{3}{l}{\textbf{Low Priority \& Low Impact (Eliminate or Delay)}} \\

Extensive UI customization & Custom themes, animations, or unnecessary aesthetics. & Keep UI minimal for now. \\

Full-scale competitor comparison & While helpful, extensive comparisons may not be crucial. & Summarize findings instead of deep-dive analysis. \\

\bottomrule
\end{tabularx}
\label{projectchoice:prioritymatrix}